% Plantilla para diapositivas de defensa oral (Beamer)
% Para DD1 (Evaluacion de Desempeno 1 - ABET 3) y todas las exposiciones
% Compilar: pdflatex slides_defensa.tex (dos pasadas)

\documentclass[aspectratio=169]{beamer}

\usepackage[utf8]{inputenc}
\usepackage[T1]{fontenc}
\usepackage[spanish]{babel}
\usepackage{graphicx}
\usepackage{booktabs}
\usepackage{hyperref}

\usetheme{Madrid}
\usecolortheme{seahorse}

% Quitar simbolos de navegacion
\setbeamertemplate{navigation symbols}{}

% --- Metadatos ---
\title{[TITULO DEL PROYECTO]}
\subtitle{[SUBTITULO - curso 1AEL0260]}
\author{Renzo Francisco Albatrino Aza}
\institute{Universidad Peruana de Ciencias Aplicadas}
\date{\today}

\begin{document}

\begin{frame}
    \titlepage
\end{frame}

\begin{frame}{Contenido}
    \tableofcontents
\end{frame}

% ============================================================
\section{Introduccion}
% ============================================================

\begin{frame}{Introduccion}
    \begin{itemize}
        \item [Contexto del problema]
        \item [Por que importa]
        \item [Contribucion propuesta]
    \end{itemize}
\end{frame}

% ============================================================
\section{Problema}
% ============================================================

\begin{frame}{Situacion problematica}
    [Descripcion del problema]
\end{frame}

\begin{frame}{Pregunta de investigacion}
    [Pregunta de investigacion]
\end{frame}

% ============================================================
\section{Estado del arte}
% ============================================================

\begin{frame}{Estado del arte}
    \begin{itemize}
        \item [Trabajo previo 1]
        \item [Trabajo previo 2]
        \item [Brecha identificada]
    \end{itemize}
\end{frame}

% ============================================================
\section{Solucion propuesta}
% ============================================================

\begin{frame}{Arquitectura}
    % \begin{figure}
    %     \centering
    %     \includegraphics[width=0.8\textwidth]{figuras/diagrama_bloques}
    % \end{figure}
    [Descripcion de la arquitectura]
\end{frame}

\begin{frame}{Objetivos}
    \begin{itemize}
        \item \textbf{Objetivo general:} [objetivo]
    \end{itemize}
    \vspace{0.5cm}
    \textbf{Objetivos especificos:}
    \begin{itemize}
        \item OE1: [objetivo 1]
        \item OE2: [objetivo 2]
        \item OE3: [objetivo 3]
    \end{itemize}
\end{frame}

% ============================================================
\section{Resultados}
% ============================================================

\begin{frame}{Resultados preliminares}
    [Resultados o esperados]
    % \begin{table}
    %     \centering
    %     \caption{Comparacion de resultados}
    %     \begin{tabular}{lcc}
    %         \toprule
    %         Metrica & Baseline & Propuesta \\
    %         \midrule
    %         Latencia & X ms & Y ms \\
    %         Energia & X mW & Y mW \\
    %         \bottomrule
    %     \end{tabular}
    % \end{table}
\end{frame}

% ============================================================
\section{Viabilidad}
% ============================================================

\begin{frame}{Viabilidad}
    \begin{itemize}
        \item \textbf{Tecnica:} [viabilidad tecnica]
        \item \textbf{Economica:} [viabilidad economica]
        \item \textbf{Social:} [viabilidad social]
        \item \textbf{Operativa:} [viabilidad operativa]
    \end{itemize}
\end{frame}

% ============================================================
\section{Conclusiones}
% ============================================================

\begin{frame}{Conclusiones}
    \begin{itemize}
        \item [Conclusion 1]
        \item [Conclusion 2]
        \item [Conclusion 3]
    \end{itemize}
\end{frame}

\begin{frame}{Limitaciones}
    \begin{itemize}
        \item [Limitacion 1]
        \item [Limitacion 2]
    \end{itemize}
\end{frame}

% ============================================================
\begin{frame}[standout]
    Gracias \\[0.5cm]
    Preguntas
\end{frame}

\end{document}
